%=======================================================================
%  Interfacial Charge Relaxation Governs the Capacitive Memory Window
%  in Ferroelectric Hf0.5Zr0.5O2 Capacitors
%
%  IEEE journal format (IEEEtran).  For IEEE Electron Device Letters use
%  \documentclass[journal]{IEEEtran} and observe the 3-page limit; for
%  IEEE TED / JEDS the present length is appropriate.
%
%  Required figure files in the same directory:
%     Fig1.png  Fig2.png  Fig3.png  Fig4.png  Fig5.png
%=======================================================================
\documentclass[journal]{IEEEtran}

\usepackage{graphicx}
\usepackage{amsmath,amssymb}
\usepackage{textcomp}
\usepackage{booktabs}
\usepackage{url}
\usepackage[caption=false,font=footnotesize]{subfig}

% ---- convenience macros -----------------------------------------------
\newcommand{\hzo}{Hf$_{0.5}$Zr$_{0.5}$O$_{2}$}
\newcommand{\xir}{\xi_{r}}
\newcommand{\CMW}{\mathrm{CMW}}
\newcommand{\fc}{f_{c}}
\newcommand{\Ecross}{E_{\mathrm{cross}}}
\newcommand{\uCcm}{\,$\mu$C/cm$^{2}$}
\newcommand{\MVcm}{\,MV/cm}

\begin{document}

\title{Interfacial Charge Relaxation Governs the Capacitive Memory Window
in Ferroelectric \hzo{} Capacitors}

\author{Apu~Das,
        and~Sourav~De,~\IEEEmembership{Senior~Member,~IEEE}
\thanks{The authors are with the College of Semiconductor Research,
National Tsing Hua University, Hsinchu 300044, Taiwan
(e-mail: apudas@nthu.edu.tw; sourav.de@mx.nthu.edu.tw).}
\thanks{Manuscript received XX XXXX; revised XX XXXX.}}



\maketitle

%=======================================================================
\begin{abstract}
The capacitive memory window (CMW) enables non-destructive readout (NDRO)
in ferroelectric capacitors, but what sets its limits has remained
unresolved. We show that although the polarization defines the two states,
the zero-bias CMW that distinguishes them is gated by an interfacial charge
layer. In TaN/\hzo{}/TiN capacitors the CMW rolls off as a single Debye
relaxation, $\fc = 179$~kHz, with Kramers--Kronig closure to 2\%; no
component survives on the ferroelectric's own timescale. At 10~K the corner
falls below 1~kHz while switching persists. Two stacks sharing the same
dielectric ceiling relax $126\times$ apart, and a series extraction shows
their interfacial layers differ by 26\% in areal capacitance but
$157\times$ in areal resistance---the read speed is set by interfacial
conductivity, not thickness. The NDRO read band is therefore bounded by an
interface-determined $\fc(T)$ varying two decades between stacks, making
single-frequency CMW values non-portable.
\end{abstract}

\begin{IEEEkeywords}
Ferroelectric capacitor, hafnium zirconium oxide, capacitive memory window,
non-destructive readout, dielectric relaxation, Maxwell--Wagner, interfacial
layer, PUND, cryogenic memory.
\end{IEEEkeywords}

\IEEEpeerreviewmaketitle

%=======================================================================
\section{Introduction}
%=======================================================================
\IEEEPARstart{F}{erroelectricity} in doped HfO$_2$~\cite{boscke2011} and
\hzo{}~\cite{muller2012} has reopened the prospect of scalable ferroelectric
memory: unlike the perovskites it displaces, \hzo{} is CMOS- and
BEOL-compatible, retains polarization in films only a few nanometres thick,
and is deposited by conformal ALD~\cite{cho2024beol,schroeder2022}.
Ferroelectric capacitors (FeCaps) and field-effect transistors (FeFETs) built
on this material are advanced candidates for embedded non-volatile memory,
storage-class memory, and in-memory and neuromorphic
computing~\cite{paul2025review,sunil2024cam,das2026cryo,sk2025stdp}.

The conventional readout of a FeCap is destructive. Sensing the switched
charge $2P_r$ requires driving the cell through its coercive field,
destroying the stored state and mandating a write-back---costing energy,
latency, and one cycle of the device's finite endurance budget. This penalty
is worst for read-dominated workloads and for in-memory computing, where the
same weights are interrogated repeatedly.

A capacitive read removes it. Because the small-signal permittivity of a
poled ferroelectric depends on which remanent state it occupies, the two
states can be distinguished by a low-amplitude, sub-coercive capacitance
measurement that leaves the polarization undisturbed. The difference between
the two state permittivities---the capacitive memory window (CMW)---is the
read quantity of NDRO FeRAM and of capacitor-based in-memory computing, where
a state-dependent capacitance directly weights accumulated
charge~\cite{mukherjee2023edl,mukherjee2023iedm,mukherjee2024iedm,lee2025aem}.
The non-volatile small-signal capacitance underlying this readout has been
documented in \hzo{} capacitors~\cite{luo2020}.

Work on this scheme has concentrated almost exclusively on enlarging the
window, and two opposing design philosophies have emerged. The first
deliberately engineers an asymmetry between the two ferroelectric/metal
interfaces so that a finite window exists at zero applied bias, reporting CMW
rising from $\approx 4.7$ to a record 8.7 with read endurance beyond
$10^{11}$ cycles~\cite{mukherjee2023edl,mukherjee2023iedm}, later enhanced
with a semiconducting IGZO interlayer and demonstrated in a 3D
architecture~\cite{mukherjee2024iedm,mukherjee2025edl}. The attraction is
circuit simplicity: a zero-bias read requires no bias generator, no
per-device read-voltage calibration, and carries no risk that the read bias
itself disturbs the state or injects leakage. The second keeps the
metal--ferroelectric--metal stack deliberately symmetric---which by
construction yields no zero-bias window---and opens the window with a
sub-coercive read bias $V_R < V_c$, on the explicit argument that built-in
asymmetry raises reliability and device-to-device variability
concerns~\cite{lee2025aem}. Both have been demonstrated at array level, and
both quote a CMW measured at a single frequency near room temperature.

What neither establishes is \emph{what governs} the window. The controlling
physics is left implicit, and with it the operating limits: usable read
bandwidth, temperature range, and stability of the read margin. This is
decisive rather than academic. If the window is limited by the ferroelectric
layer it inherits that layer's
robustness~\cite{pesic2016,schenk2014,grimley2016}; if it is limited by
interfacial charge, its bandwidth and temperature range are governed by
space-charge and trap dynamics, which are neither as fast nor as stable and
which vary strongly with electrode and interfacial chemistry.

A symmetry argument shows the polarization alone cannot be the answer. The
polarization is \emph{necessary}, since it defines the two states; without
ferroelectric switching there is no stored bit. It is not \emph{sufficient}:
for an ideally symmetric capacitor the small-signal permittivity is an even
function of the polarization state, $\xi(+P_r) = \xi(-P_r)$, so the two
$C$--$V$ branches intersect at zero bias and $\CMW \equiv 0$
\emph{regardless of the magnitude of} $2P_r$. A finite zero-bias window does
not follow from having a large polarization---it requires an asymmetry
between the two states. This is implicit in the need to engineer interfacial
asymmetry to open a window~\cite{mukherjee2023edl,mukherjee2023iedm}, and in
the absence of a zero-bias window in symmetric stacks~\cite{lee2025aem}. The
questions addressed here are what that asymmetry is, and what governs the
window's dynamics.

There is strong independent reason to suspect the interfacial layer. In \hzo{}
capacitors, an oxygen-deficient interfacial region forms unavoidably against
nitride electrodes: TiN and TaN scavenge oxygen from the film during
deposition and annealing, forming TiO$_x$N$_y$ or TaO$_x$N$_y$ and generating
oxygen vacancies in the adjacent
\hzo{}~\cite{kim2023vacancy,segantini2023strain,koroleva2022}. Systematic
comparison across nitride, metal and oxide electrodes identifies this
electrode--ferroelectric interface as the \emph{primary} constraint on
endurance and retention, with impedance spectroscopy of 10~nm \hzo{}
capacitors resolving distinct interface and bulk RC elements whose
resistances differ by more than five orders of
magnitude~\cite{alcala2023}. The same interfacial region drives the wake-up
transformation and the accompanying
imprint~\cite{pesic2016,schenk2014,grimley2016,schenk2014apr,li2024imprint},
and interfacial-layer soft breakdown has been proposed as its microscopic
mechanism in ultrathin \hzo{}~\cite{cho2024wakeup,de2024tradeoff}. Vacancy
distribution and phase composition in scaled \hzo{} capacitors have been
mapped directly~\cite{iung2025,kumar2025review}, and charge-pumping analyses
show interface-trap accumulation, rather than intrinsic ferroelectric
fatigue, dominates endurance degradation in HfO$_2$
FeFETs~\cite{rana2025cp,raffel2025,senapati2025}. Most directly, examination
of the same electrically degraded gate stack after memory-window closure
shows robust polarization switching is preserved, establishing that the loss
of device functionality is
interface-driven~\cite{das2026paradox,das2025sisc}. Interface engineering is
correspondingly the principal lever on FeCap reliability: the engineered
WO$_x$/\hzo{}/PEALD-TiN stack studied here as Stack~B achieves endurance
beyond $2\times10^{10}$ cycles~\cite{paul2026wox}.

If the interfacial layer governs wake-up, imprint and endurance, it is
natural to ask whether it also governs the capacitive read. The two candidate
pictures make opposing, testable predictions. If the window is
\emph{ferroelectric-limited}---arising from the field-dependent dielectric
response of the ferroelectric itself, as in Landau or domain-wall treatments
of the butterfly $C$--$V$~\cite{lee2025aem,tagantsev2002,segantini2022ti}---it
should relax on the ferroelectric's own timescale, which persists into the
GHz range~\cite{kim2023impedance}, be insensitive to the electrode stack, and
show the broadening characteristic of a distribution of pinning
energies~\cite{jonscher1977,jonscher1983}. If it is
\emph{interface-limited}, it should exhibit a single relaxation whose time
constant is set by the interfacial layer alone, change with the stack, and
freeze on cooling while the polarization remains operational.

We resolve which holds using frequency-resolved $C$--$V$ spectroscopy from
1~kHz to 4~MHz on two stacks, with PUND and switching-current measurements at
300~K and 10~K.

%=======================================================================
\section{Device Fabrication and Measurement}
%=======================================================================

\subsection{Device Stacks}
Two metal--ferroelectric--metal capacitors were compared
[Fig.~\ref{fig:stacks}(a)]. \textbf{Stack~A} is
TaN(120~nm)/\hzo{}(10~nm)/TiN(50~nm) on Si. \textbf{Stack~B} is
PE-TiN/TiN/\hzo{}/WO$_x$/W on Si, incorporating a tungsten-oxide interfacial
layer between the ferroelectric and the bottom electrode~\cite{paul2026wox}.
\hzo{} was deposited by ALD and crystallized into the polar orthorhombic
phase by rapid thermal annealing with the top electrode present, providing
the mechanical constraint required for metastable phase stabilization, at a
BEOL-compatible thermal budget ($<400\,^\circ$C)~\cite{cho2024beol}. Phase
composition was verified by grazing-incidence XRD [Fig.~\ref{fig:stacks}(b)].

Neither stack contains a deliberately deposited interfacial layer at its
nitride contacts. However, an oxygen-deficient interfacial region forms
unavoidably at nitride/\hzo{} interfaces during deposition and
crystallization, as TiN and TaN scavenge oxygen from the
film~\cite{kim2023vacancy,segantini2023strain,koroleva2022,alcala2023}. In
Stack~A the two electrodes are chemically dissimilar (TaN and TiN), so the two
interfacial regions differ, providing the asymmetry required for a finite
zero-bias window; in Stack~B the asymmetry is set deliberately by the WO$_x$
layer~\cite{paul2026wox}.

\subsection{Small-Signal $C$--$V$ Spectroscopy}
$C$--$V$ characteristics were measured in parallel ($C_p$--$R_p$) mode using
an LCR meter with a 50~mV rms small-signal amplitude, at 61 logarithmically
spaced frequencies from 1~kHz to 4~MHz. The DC bias was swept
$-2$~V~$\rightarrow +2$~V~$\rightarrow -2$~V to traverse both polarization
branches in a single acquisition, so that the branch reached after negative
poling ($\xi_\uparrow$) and after positive poling ($\xi_\downarrow$) are
recorded under identical conditions.

The effective relative permittivity was obtained from the measured parallel
capacitance as
\begin{equation}
\xir(V,f) \;=\; \frac{C_p(V,f)\,d}{\varepsilon_0 A} \;=\; \frac{C_p(V,f)}{C_0}
\label{eq:eps}
\end{equation}
with $d$ the \hzo{} thickness, $A$ the capacitor area, and
$C_0 = \varepsilon_0 A/d$ the geometric capacitance. For Stack~A,
$C_0 = 15.0$~pF.

For each frequency the two branches were separated by the sign of $dV/dt$ and
interpolated to zero bias, giving $\xi_\uparrow(0)$ and $\xi_\downarrow(0)$.
The capacitive memory window and the crossing field are defined as
\begin{equation}
\CMW_\xi(f) = \xi_\downarrow(0,f) - \xi_\uparrow(0,f),
\qquad
\xi_\downarrow(\Ecross) = \xi_\uparrow(\Ecross)
\label{eq:cmw}
\end{equation}
$\Ecross$ being the field at which the two branches intersect and the states
become capacitively indistinguishable.

The differential dielectric loss was obtained from the measured branch
conductances $G = 1/R_p$, again referenced to the geometric capacitance:
\begin{equation}
\Delta\xi''(f) \;=\; \frac{G_\downarrow(0,f) - G_\uparrow(0,f)}{\omega\,C_0},
\qquad \omega = 2\pi f
\label{eq:loss}
\end{equation}
Because \eqref{eq:loss} uses the conductance channel while \eqref{eq:cmw}
uses the capacitance channel, the real and imaginary parts of the
differential response are determined independently, permitting a
Kramers--Kronig consistency test~\cite{nicollian1982,schroder2006}. The same
conductance-based approach is used to quantify ferroelectric/dielectric
interface traps~\cite{li2021traps}.

\subsection{Series-Resistance De-Embedding}
Above $\sim$1~MHz the measured response is dominated by the fixture and
electrode resistance rather than by the dielectric. The apparent parallel
capacitance of a series $R_s$--$L$--$C$ network is
\begin{equation}
C_{p,\mathrm{meas}}(f) \;=\; \frac{1}{\omega}\,
\mathrm{Im}\!\left[\frac{1}{R_s + j\omega L + (j\omega C)^{-1}}\right]
\label{eq:rlc}
\end{equation}
Fitting \eqref{eq:rlc} to the 0.8--4~MHz data of Stack~A gives
$R_s = 267\,\Omega$, $L = 14\,\mu$H and $C = 329$~pF, placing the RC corner
at $f_{RC} = 1/(2\pi R_s C) = 1.81$~MHz and the LC resonance at
$f_0 = 1/(2\pi\sqrt{LC}) = 2.34$~MHz, against a measured zero crossing of the
apparent permittivity at 2.44~MHz. The fractional error in $C_p$ from series
resistance is $(\omega R_s C)^2$, which equals 0.010 at $\fc$ and 0.076 at
500~kHz; all main-text analysis is therefore restricted to $f \le 500$~kHz.

\subsection{PUND and Switching-Current Measurements}
Polarization was measured by the positive-up-negative-down (PUND) sequence,
in which a preset pulse is followed by two positive pulses (P, U) and two
negative pulses (N, D). The first pulse of each polarity switches the
polarization and the second does not, so subtracting the two removes the
non-switching dielectric and leakage contributions:
\begin{equation}
2P_r \;=\; \frac{1}{A}\int \left[I_P(t) - I_U(t)\right]\,dt
\label{eq:pund}
\end{equation}
Triangular sweeps with a half-period $T_{pw} = 500\,\mu$s were used for the
switching-current characteristics of Fig.~\ref{fig:temperature}(a); the
switching current is the difference between the switching and non-switching
branches, $I_{sw} = I_P - I_U$. Coercive voltages were extracted from the
zero crossings of $P(V)$ on each branch, and the loop imprint from their
midpoint:
\begin{equation}
\begin{aligned}
&P(V_c^{\pm}) = 0, \qquad
V_{\mathrm{imp}} = \tfrac{1}{2}\left(V_c^{+} + V_c^{-}\right),\\[2pt]
&E_c = \tfrac{1}{2}\left(|V_c^{+}| + |V_c^{-}|\right)\big/ d
\end{aligned}
\label{eq:coercive}
\end{equation}
Switching kinetics were characterized by a pulse-width series, applying
rectangular write pulses of amplitude $V_R = \pm2$~V and width
10~$\mu$s--1~ms and reading $2P_r$ by PUND after each write.

\subsection{Cryogenic Measurement}
Low-temperature measurements were performed in a closed-cycle helium
cryogenic probe station. The sample stage was cooled by a Gifford--McMahon
helium compressor to a base temperature of 10~K under vacuum, with the
temperature stabilized to within $\pm0.1$~K and a dwell of at least 30~min
before each measurement to ensure thermal equilibrium of the probe arms.
Identical probes, cabling and instrument ranges were used at 300~K and 10~K
so that the two datasets are directly comparable; in particular, the series
resistance obtained above applies to both.

\subsection{Model Fitting}
The window and its loss were fitted to the two-component Debye form of
Section~\ref{sec:debye} by non-linear least squares. Fit windows were
$f \le 600$~kHz for Stack~A and $f \le 20$~kHz for Stack~B for the CMW, and
$f \le 300$~kHz and $f \le 30$~kHz respectively for the peak permittivity,
chosen to exclude the series-resistance-affected region. The Cole--Cole locus
was fitted over 30--300~kHz by a circular arc constrained to pass through the
origin, since $\Delta\xi'$ and $\Delta\xi''$ both vanish as
$f \rightarrow \infty$; the depression angle $\alpha$ and the Cole--Cole
exponent are related by $\beta = 1 - 2\alpha/\pi$.

%=======================================================================
\section{Results and Discussion}
%=======================================================================

\subsection{The Capacitive Memory Window and Its Frequency Collapse}
GIXRD [Fig.~\ref{fig:stacks}(b)] shows the dominant polar orthorhombic
reflection near $2\theta \approx 30.4^\circ$ with a residual monoclinic
contribution near $28.5^\circ$, confirming both films as ferroelectric
\hzo{}~\cite{boscke2011,muller2012,iung2025}. Fig.~\ref{fig:stacks}(c)
summarizes the wake-up transformation for
reference~\cite{pesic2016,schenk2014,grimley2016,cho2024wakeup}.

Fig.~\ref{fig:cmwdef}(a) shows the effective permittivity of Stack~A at
1~kHz for both sweep directions. The butterfly peaks reach
$\xir \approx 62$ near $\mp0.55$\MVcm{}, and the two branches do not coincide
at zero field: $\xi_\downarrow(0) = 55.0$ and $\xi_\uparrow(0) = 49.8$,
giving $\CMW_\xi \approx 5.2$. This splitting---read at zero bias without
disturbing the stored state---is the capacitive memory window. Stack~B
[Fig.~\ref{fig:cmwdef}(b)] reaches the same peak permittivity,
$\xir \approx 62$, but gives $\CMW_\xi \approx 2.3$ at the same 1~kHz probe
frequency.

Fig.~\ref{fig:cmwdef}(c) follows Stack~A from 1~kHz to 400~kHz. Two things
occur together: the peaks fall from 62 to 30, and the branch separation
shrinks until, by 400~kHz, the window has essentially closed. Critically, the
dielectric response itself does not vanish---a permittivity of $\approx$30
persists at 400~kHz and $\approx$27 at 500~kHz. What disappears is the
\emph{difference} between the states. The window is therefore carried by a
dispersive contribution superposed on a largely non-dispersive background,
motivating the two-component treatment below.

\subsection{The Window Relaxes as a Single Debye Species}
\label{sec:debye}
We model the poled stack as a two-component dielectric,
\begin{equation}
\xi^{*}_{\pm}(\omega) \;=\; \xi_\infty + \frac{\Delta\xi_\pm}{1+j\omega\tau}
\label{eq:twocomp}
\end{equation}
where $\xi_\infty$ is the polarization-independent high-frequency
permittivity and $\Delta\xi_\pm$ the state-dependent relaxing term. Because
$\xi_\infty$ cancels analytically in the difference between states, the
measured window isolates the relaxing component and is insensitive to the
large non-relaxing background:
\begin{align}
\CMW(f) &= \frac{\CMW_0}{1+(f/\fc)^2} \label{eq:debyereal}\\
\Delta\xi''(f) &= \frac{A_{\mathrm{DC}}}{f}
   + \CMW_0\,\frac{f/\fc}{1+(f/\fc)^2} \label{eq:debyeimag}
\end{align}
with $\fc = 1/2\pi\tau$. The first loss term accounts for a state-dependent
DC leakage contribution~\cite{nicollian1982,schroder2006}.

\textit{Real part.} Fig.~\ref{fig:relaxation}(a) shows the measured CMW for
both stacks across 61 frequencies. Stack~A is described by
\eqref{eq:debyereal} with $\CMW_0 = 5.38$ and $\fc = 179$~kHz, corresponding
to $\tau = 0.90\,\mu$s. Stack~B gives $\CMW_0 = 3.6$ and $\fc = 1.4$~kHz,
i.e.\ $\tau = 114\,\mu$s; because Stack~B's roll-off begins below 10~kHz its
plateau is an extrapolation, and its measured 1~kHz value (2.36) already lies
at $f/\fc = 0.71$ on the roll-off.

\textit{Imaginary part.} Fig.~\ref{fig:relaxation}(b) fits
\eqref{eq:debyeimag} to the differential loss of Stack~A, obtained
independently from the conductance channel. The decomposition returns a
leakage coefficient $A_{\mathrm{DC}} = 551$~Hz and a Debye peak at
$\fc = 182$~kHz with amplitude 5.94---within 2\% of the real-part corner. The
residual is 0.13 against a peak loss of 2.9.

\textit{Cole--Cole locus.} Fig.~\ref{fig:relaxation}(c) presents the
differential Cole--Cole plot over 30--300~kHz after removal of the leakage
term. Below 30~kHz the subtraction dominates the signal and above 300~kHz
series-resistance loss lifts the locus; those points are shown as open
symbols and excluded from the fit. The circular-arc fit constrained to pass
through the origin returns a depression angle $\alpha = -2.7^\circ$, zero
within uncertainty and of unphysical sign, so $\beta = 1$ within resolution.
Forcing an ideal Debye arc with its centre on the real axis increases the rms
radial residual only from 3.1\% to 3.5\% of the radius, so a distribution of
relaxation times is not required by the data. The arc amplitude, 5.50, agrees
to 2.1\% with $\CMW_0 = 5.38$ from the real-part fit---Kramers--Kronig
closure between two independent measurement channels~\cite{colecole1941}.

Three determinations of the corner---179~kHz (real part), 182~kHz (loss peak)
and the arc geometry---agree within 2\%. The window relaxes as one
well-defined species, which a distributed domain-wall or defect-continuum
response would not produce~\cite{jonscher1977,jonscher1983}.

\textit{Completeness of the relaxation.} Admitting a non-relaxing offset,
$\CMW(f) = \CMW_\infty + \CMW_0/[1+(f/\fc)^2]$, returns
$\CMW_\infty = -0.36 \pm 0.11$, negative and therefore unphysical as a
residual; any non-relaxing component is below 7\% of $\CMW_0$. This is
significant because a window limited by the ferroelectric's own response,
which persists into the GHz range~\cite{kim2023impedance}, would leave such a
residual. Instead, at 500~kHz the window has fallen to 0.45 while the
permittivity is still 27: what survives the relaxation is state-independent.

\textit{Intrinsic origin of the corner.} With $R_s = 267\,\Omega$, the
circuit-limited corner lies at 1.81~MHz and the correction factor
$(\omega R_s C)^2 = 0.010$ at $\fc$. The 179~kHz corner is therefore a
property of the device, not the measurement circuit.

\subsection{Temperature: The Window Freezes While Switching Persists}
Fig.~\ref{fig:temperature} tests whether the window follows the polarization
on the temperature axis.

\textit{Switching survives.} Fig.~\ref{fig:temperature}(a) shows the
switching current at both temperatures for $T_{pw} = 500\,\mu$s. At 300~K the
peaks reach $\pm65\,\mu$A at $\mp0.99$\MVcm{}. At 10~K the peaks fall to
$\approx13\,\mu$A and broaden, displaced to higher field, but switching
persists. Fig.~\ref{fig:temperature}(b) shows the corresponding $P$--$E$
loops: $2P_r$ falls from 33.1 to 20.3\uCcm{} while the coercive field rises
from $E_c = 0.93$ to 1.17\MVcm{}, a 26\% increase.

\textit{The reduction is kinetic, not a loss of ferroelectricity.}
Fig.~\ref{fig:temperature}(c) resolves this directly. At 300~K, $2P_r$ rises
steeply with write pulse width and saturates at 33\uCcm{} by 200~$\mu$s. At
10~K it is still increasing at 1~ms, reaching only 20\uCcm{}, and has not
saturated within the accessible pulse range. At fixed drive amplitude and
pulse width the film is therefore progressively under-driven on cooling,
consistent with the reduced domain-nucleation rate expected from NLS
kinetics~\cite{tagantsev2002} and with reports of slowed switching in scaled
\hzo{} at cryogenic temperature~\cite{das2026cryo}. The ferroelectric remains
fully functional at 10~K---it is simply slower.

\textit{The window behaves oppositely.} At 1~kHz the window falls from 5.2 to
2.2 on cooling---a 58\% reduction but not extinction---and above 1.1~kHz
drops to the noise floor of $\approx$0.2 [Fig.~\ref{fig:temperature}(d)].
Since no Debye function can fall tenfold over a 10\% frequency interval, this
is not a resolved roll-off but a bound: $\fc \lesssim 1$~kHz at 10~K, more
than two decades below its room-temperature value. Notably the window does
not vanish at 10~K; its \emph{clock} does, so that at sufficiently low
frequency the interfacial charge can still follow.

This is the first indication that the window is interface-limited: at 10~K
the polarization still switches while the capacitive readout has lost its
bandwidth. A quantity limited by the ferroelectric would not lose bandwidth
while the ferroelectric itself remains operational.

\textit{The asymmetry is not the ferroelectric imprint.}
Fig.~\ref{fig:temperature}(e) compares the two independently measured
built-in voltages. From the $P$--$E$ loops, the imprint \eqref{eq:coercive}
is $+0.017$~V at 300~K and $+0.009$~V at 10~K---essentially symmetric, since
$V_c^{+} = 0.949$~V and $V_c^{-} = -0.915$~V at 300~K. The $C$--$V$ crossing
voltage of the same stack is $0.221 \pm 0.005$~V, thirteen times larger, and
is frequency-independent: $0.219 \pm 0.003$~V over 1--10~kHz,
$0.224 \pm 0.005$~V over 10--100~kHz and 0.215~V over 100--300~kHz, i.e.\
constant to $\pm2$\% across 2.5 decades while the window amplitude falls from
5.4 to below 1. The asymmetry gating the window is therefore not the
ferroelectric loop asymmetry. It is what would be expected of a thin
interfacial layer supporting a significant fraction of the applied bias while
contributing negligibly to the switched
charge~\cite{alcala2023,schenk2014apr,li2024imprint}. The dichotomy---a
relaxing amplitude on a static position---is itself difficult to reconcile
with a mechanism intrinsic to the ferroelectric, since a field-dependent
ferroelectric process would be expected to shift the crossing point as it
disperses.

\subsection{The Interface Sets the Clock}
If the window is interface-limited, changing the interface should change its
dynamics; if ferroelectric-limited, it should not.

\textit{Identical ferroelectric ceilings.} Fig.~\ref{fig:twoclocks}(a) shows
the peak permittivity of both stacks. Extrapolating the roll-off to zero
frequency using the same Debye form gives $\xir(f\!\rightarrow\!0) =
62.4 \pm 1.2$ for Stack~A and $62.7 \pm 0.9$ for Stack~B---statistically
identical, and consistent with the measured 1~kHz values of 62.5 and 62.2.
Both stacks therefore present the same intrinsic dielectric amplitude from
the \hzo{} layer. Their high-frequency limits differ,
$\xi_\infty = 21.7 \pm 0.9$ and $15.8 \pm 0.5$, as expected if the series
interfacial layers differ. (Stack~B's apparent decline to $\xir \approx 6$
above 200~kHz falls below its own fitted $\xi_\infty$ and is circuit-limited,
its larger capacitance placing its RC corner lower.)

\textit{Two-decade difference in clock.} Against those identical ceilings,
the two windows relax $126\times$ apart: $\tau = 0.90\,\mu$s versus
114~$\mu$s. Fig.~\ref{fig:twoclocks}(b) shows that scaling each dataset by
its own corner frequency and plateau collapses both onto the universal Debye
form $1/(1+x^2)$---the same functional response, two different clocks.
Fig.~\ref{fig:twoclocks}(c) summarizes: identical dielectric response (62
versus 62), windows of 5.2 versus 2.4 at 1~kHz, and built-in fields of 0.22
versus 0.08\MVcm{}. The stack with the larger built-in field also shows the
larger window, consistent with the asymmetry setting the window amplitude,
although two devices cannot establish a scaling relation between them.

\textit{Controls.} First, Stack~B's low corner is not an RC artifact: with
the fitted $R_s$ and Stack~B's capacitance the circuit-limited corner lies
near 220~kHz, more than two decades above the observed 1.4~kHz. Second, the
difference cannot originate in the ferroelectric, whose low-frequency
dielectric ceiling is identical in both to within 1.5\%.

We attribute the difference to the interface, noting the structural
correlation that Stack~B, which contains a deliberate WO$_x$ interfacial
layer~\cite{paul2026wox}, is the slower of the two---as expected for
relaxation across a more resistive interfacial layer. We state the limitation
explicitly: the two stacks differ in electrode materials and ferroelectric
thickness as well as in the WO$_x$ layer, so a single causal variable cannot
be isolated from two devices. The claim the data support unambiguously is
that two stacks sharing the same ferroelectric dielectric ceiling exhibit
windows differing by two orders of magnitude in relaxation time---the clock
is not set by the ferroelectric.

\subsection{Mechanism: Maxwell--Wagner Relaxation}
A ferroelectric layer in series with a thinner, more resistive interfacial
layer exhibits Maxwell--Wagner interfacial relaxation with a single time
constant~\cite{jonscher1983,kremer2003}
\begin{equation}
\tau = \frac{\varepsilon_i d_f + \varepsilon_f d_i}
            {\sigma_i d_f + \sigma_f d_i}
     \;\approx\; \frac{\varepsilon_i}{\sigma_i}
     \qquad (d_i \ll d_f)
\label{eq:mw}
\end{equation}
set by the interfacial layer alone [Fig.~\ref{fig:mechanism}(a),(b)]. This
accounts for every observation without additional assumptions.

A two-layer series network produces exactly one relaxation, matching the
undepressed arc of Fig.~\ref{fig:relaxation}(c) and the 2\% agreement among
three independent corner determinations. The interfacial layer supports a
static built-in potential set by the electrode/oxide chemistry, giving the
frequency-independent $\Ecross$ above while the charge redistribution across
it relaxes. Because $\tau$ depends only on the interfacial layer, it changes
with the stack: taking $\varepsilon_i$ from the measured $\xi_\infty$ values,
\eqref{eq:mw} implies $\sigma_i = 2.2\times10^{-4}$~S/m for Stack~A and
$1.2\times10^{-6}$~S/m for Stack~B. These are order-of-magnitude estimates,
but they lie in the range expected for a defective, vacancy-rich interfacial
oxide~\cite{kim2023vacancy,segantini2023strain,koroleva2022,iung2025}, and
their ratio reproduces the observed spread in $\tau$. The picture is
independently supported by impedance spectroscopy of 10~nm \hzo{} capacitors,
which resolves distinct interface and bulk RC elements with pristine
resistances of $\approx$1~k$\Omega$ and $\approx$500~M$\Omega$
respectively~\cite{alcala2023}. Since $\sigma_i$ is thermally activated,
$\tau$ grows rapidly on cooling and the corner falls out of the measurement
band [Fig.~\ref{fig:temperature}(d)], while the ferroelectric---whose
switching does not depend on interfacial conduction---continues to operate.

\textit{Separating the interfacial capacitance from the interfacial
resistance.} Because the interfacial region enters as a series element, its
areal capacitance can be extracted directly from the two permittivity limits
without any assumption about its thickness or permittivity. Writing the
low-frequency limit $\xi_{FE}$ (interfacial layer conducting) and the
high-frequency limit $\xi_\infty$ (interfacial layer blocking) as series
capacitances gives
\begin{equation}
c_{IL} = \frac{\varepsilon_0}{d\left(\xi_\infty^{-1}-\xi_{FE}^{-1}\right)},
\qquad
R_{IL} = \frac{\tau}{c_{IL}}
\label{eq:cil}
\end{equation}
the second relation following from $\tau = R_{IL}c_{IL}$ for the series
element. Applying \eqref{eq:cil} to the fitted limits gives
$c_{IL} = 2.95$ and 2.34~$\mu$F/cm$^2$ for Stacks~A and~B---a difference of
only 26\%---while the corresponding areal resistances are 0.31 and
48.8~$\Omega\cdot$cm$^2$, a factor of 157. The product reproduces the measured
relaxation times to within 3\% ($157 \times 0.79 = 124$ versus the observed
$126\times$), confirming internal consistency. The two stacks therefore
present interfacial layers of nearly the same series capacitance but
differing by more than two orders of magnitude in resistance: the WO$_x$
layer acts by lowering interfacial conductance, not by adding series
thickness [Fig.~\ref{fig:mechanism}(c)]. This is a design statement---to move the read corner, engineer
interfacial conductivity rather than interfacial thickness---and it is
obtained without imaging the interface. We note that $c_{IL}$ and $R_{IL}$
are sums over both interfaces of each stack and cannot be apportioned between
them electrically; converting them to a physical thickness additionally
requires $\xi_{IL}$, for which $\sigma_i \approx 10^{-6}$--$10^{-4}$~S/m
follows if $\xi_{IL} \approx 6$--30 is assumed.

This does not imply the ferroelectric is absent from the window: below $\fc$
the measured capacitance is that of the ferroelectric layer, and the
polarization defines the two states being compared. It implies that the
window's \emph{existence} requires the interfacial asymmetry, and that its
magnitude, clock and temperature range are interfacial properties. We present
this as the interpretation consistent with all measurements; direct
structural confirmation by TEM/EELS of the two
interfaces~\cite{iung2025,kumar2025review,das2026paradox} is the natural next
step.

\subsection{Design Rules for Capacitive Readout}
\textit{The usable read band is $f \lesssim \fc(T)$, and it is an interface
property.} Reading above the corner does not merely reduce the margin, it
eliminates it. Since $\fc$ varies from 179~kHz to 1.4~kHz between our two
stacks, a read frequency well inside the plateau for one stack lies two
decades beyond the corner for the other. A CMW quoted at a single frequency
is therefore not portable across stacks, and apparent disagreements between
reports may reflect only where each probe frequency sat relative to that
device's corner. We suggest CMW values be reported together with the measured
$\fc$, or with evidence that the read frequency lies in the plateau.

\textit{A larger zero-bias window trades against stability.} The margin
obtained through interfacial asymmetry is built from the same charge that
relaxes with frequency and freezes on cooling. This gives a physical basis
for the reliability concerns motivating symmetric-stack, finite-read-bias
designs~\cite{lee2025aem,rana2025cp,fehlings2025}: the two philosophies
represent a genuine trade between circuit simplicity and read-margin
stability.

\textit{Capacitive read is a warm-temperature technique.} At 10~K the readout
bandwidth has collapsed below 1~kHz while switching persists, so a
capacitive-read cell that functions at 300~K may fail to resolve its states
at cryogenic temperature. Cryogenic ferroelectric memory should therefore
fall back on the polarization read, which remains available precisely because
the polarization survives---with write pulses lengthened to $\gtrsim$1~ms to
accommodate the slower kinetics of
Fig.~\ref{fig:temperature}(c)~\cite{das2026cryo}.

%=======================================================================
\section{Conclusion}
%=======================================================================
The capacitive memory window of ferroelectric \hzo{} capacitors is gated by
an interfacial charge layer. The polarization defines the two states, but the
interface determines whether, how fast, and over what temperature range they
can be distinguished capacitively. Symmetry requires the window to vanish for
a symmetric stack regardless of $2P_r$, and we find the interfacial asymmetry
that opens it ($\Ecross = 0.221$\MVcm{}) is an order of magnitude larger than
the ferroelectric loop imprint (0.017~V) and frequency-independent to
$\pm2$\% across 2.5 decades. The window relaxes as a single Debye species
with $\fc = 179$~kHz, confirmed by three independent determinations agreeing
within 2\% and by Kramers--Kronig closure to 2.1\%, and it relaxes
completely, leaving no component on the ferroelectric's own timescale. Its
corner falls below 1~kHz at 10~K while switching persists, the reduced
$2P_r$ being kinetically limited rather than lost. Two stacks whose
ferroelectric dielectric ceilings coincide to within 1.5\% exhibit windows
differing $126\times$ in relaxation time. The behaviour is quantitatively
consistent with Maxwell--Wagner relaxation across an interfacial layer of
implied conductivity $10^{-6}$--$10^{-4}$~S/m, thermally activated and set by
the interface chemistry. The NDRO read band is therefore bounded by an
interface-determined corner frequency varying by orders of magnitude between
stacks, single-frequency CMW values are not comparable across reports, and a
larger zero-bias window is purchased from the least stable part of the
device. Engineering the capacitive read is an exercise in interface
engineering. The evolution of the window under extended field cycling is the
subject of ongoing work.

%=======================================================================
%  FIGURES
%=======================================================================
\begin{figure*}[!t]
\centering
\includegraphics[width=\textwidth]{Fig1.png}
\caption{As-fabricated device stacks and phase characterization.
(a) Schematic of the two metal--ferroelectric--metal capacitors: Stack~A,
TaN/\hzo{}(10~nm)/TiN, and Stack~B, TiN/PE-TiN/\hzo{}(8~nm)/WO$_x$/W, the
latter incorporating a deposited tungsten-oxide layer at the bottom contact.
Layers are drawn as deposited; reacted interfacial regions are addressed in
Section~III-E. (b) GIXRD of the \hzo{} film showing the dominant polar
orthorhombic reflection with residual monoclinic contribution (inset).
(c) Schematic of the wake-up process, in which field cycling converts
non-polar monoclinic to polar orthorhombic phase and de-pins domains, with
representative pristine and woken $P$--$E$ loops.}
\label{fig:stacks}
\end{figure*}

\begin{figure*}[!t]
\centering
\includegraphics[width=\textwidth]{Fig2.png}
\caption{Definition and frequency collapse of the CMW at 300~K.
(a) Effective relative permittivity of Stack~A at 1~kHz for both sweep
directions; the branches are displaced at zero field,
$\xi_\downarrow(0) = 55.0$ and $\xi_\uparrow(0) = 49.8$, defining
$\CMW_\xi \approx 5.2$. (b) The same measurement on Stack~B, giving
$\CMW_\xi \approx 2.3$ at the same peak permittivity. (c) Stack~A from 1~kHz
to 400~kHz; the branch separation collapses while a permittivity of
$\approx$30 remains at 400~kHz, indicating that the window is carried by a
dispersive contribution superposed on a largely non-dispersive background.}
\label{fig:cmwdef}
\end{figure*}

\begin{figure*}[!t]
\centering
\includegraphics[width=\textwidth]{Fig3.png}
\caption{The window relaxes as a single Debye species. (a) CMW at zero bias
versus frequency for both stacks with single-Debye fits: $\fc = 179$~kHz
($\tau = 0.90\,\mu$s) for Stack~A and $\fc = 1.4$~kHz ($\tau = 114\,\mu$s)
for Stack~B. (b) Differential loss $\Delta\xi''$ for Stack~A, decomposed into
a DC-leakage contribution $\propto 1/f$ and a Debye loss peak at 182~kHz,
obtained from the conductance channel and therefore an independent
determination. (c) Differential Cole--Cole locus over 30--300~kHz after
removal of the leakage term, fitted by a circular arc constrained to pass
through the origin. The depression angle, $-2.7^\circ$, is zero within
uncertainty; forcing an ideal Debye arc increases the rms radial residual only
from 3.1\% to 3.5\%, and the arc amplitude, 5.50, agrees to 2.1\% with
$\CMW_0 = 5.38$ from the real-part fit (dashed). Open symbols lie outside the
analysis band.}
\label{fig:relaxation}
\end{figure*}

\begin{figure*}[!t]
\centering
\includegraphics[width=\textwidth]{Fig4.png}
\caption{The window freezes while switching persists. (a) Switching current at
300~K and 10~K ($T_{pw} = 500\,\mu$s); peaks fall from $\pm65\,\mu$A to
$\approx13\,\mu$A and shift to higher field, but switching persists.
(b) $P$--$E$ hysteresis; $2P_r$ falls from 33.1 to 20.3\uCcm{} while $E_c$
rises from 0.93 to 1.17\MVcm{}. (c) $2P_r$ versus write pulse width at
$V_R = \pm2$~V; at 300~K the polarization saturates at 33\uCcm{} by
200~$\mu$s, whereas at 10~K it has not saturated even at 1~ms, showing the
reduction to be kinetically limited. (d) CMW at 10~K, resolvable only at the
lowest measured frequency and at the noise floor above 1.1~kHz, bounding
$\fc \lesssim 1$~kHz. (e) Built-in voltage comparison: the $C$--$V$ crossing
voltage (0.221~V) exceeds the $P$--$E$ loop imprint (0.017~V at 300~K,
0.009~V at 10~K) by an order of magnitude, indicating that the asymmetry
gating the window is not the ferroelectric loop asymmetry.}
\label{fig:temperature}
\end{figure*}

\begin{figure*}[!t]
\centering
\includegraphics[width=\textwidth]{Fig5.png}
\caption{The interface, not the ferroelectric, sets the relaxation time.
(a) Peak permittivity versus frequency; both stacks extrapolate to the same
dielectric ceiling, $\xir(f\!\rightarrow\!0) = 62.4 \pm 1.2$ and
$62.7 \pm 0.9$, despite the two-decade difference in corner frequency, while
their high-frequency limits differ ($\xi_\infty = 21.7$ and 15.8).
(b) Master-curve representation: scaling each dataset by its own corner
frequency and plateau collapses both onto $1/(1+x^2)$. (c) Summary comparison
of dielectric response, window magnitude and built-in field. The stack
containing an explicit WO$_x$ interfacial layer exhibits the slower
relaxation, as expected for Maxwell--Wagner relaxation across a more
resistive interfacial layer.}
\label{fig:twoclocks}
\end{figure*}

\begin{figure*}[!t]
\centering
\includegraphics[width=\textwidth]{Fig6.png}
\caption{Interfacial origin of the capacitive memory window.
(a) Cross-sections of both stacks including the reacted interfacial regions:
TaO$_x$N$_y$ and TiO$_x$N$_y$ form at the nitride contacts of Stack~A and at
the PE-TiN contact of Stack~B by oxygen scavenging during deposition and
crystallization (dashed outlines: inferred from the electrical extraction and
from the literature, not imaged here), while the WO$_x$ layer of Stack~B is
deposited (solid outline). Markers indicate oxygen vacancies and trapped
charge in the \hzo{} adjacent to each interface. (b) Equivalent circuit: the
ferroelectric layer in series with the interfacial layer, giving a single
Maxwell--Wagner relaxation $\tau = R_{IL}c_{IL} \approx \varepsilon_i/\sigma_i$.
(c) Interfacial areal capacitance and resistance extracted using
\eqref{eq:cil}: the two stacks differ by 26\% in $c_{IL}$ but by $157\times$
in $R_{IL}$.}
\label{fig:mechanism}
\end{figure*}

%=======================================================================
%  TABLE
%=======================================================================
\begin{table}[!t]
\renewcommand{\arraystretch}{1.15}
\caption{Extracted Parameters for Both Stacks}
\label{tab:params}
\centering
\scriptsize
\setlength{\tabcolsep}{3pt}
\begin{tabular}{@{}lcc@{}}
\toprule
Parameter & Stack A & Stack B \\
\midrule
Top / bottom electrode          & TaN / TiN        & TiN / WO$_x$--W \\
Thickness $d$ (nm)              & 10               & 8$^{\dagger}$ \\
$C_0$ (pF)                      & 15.0             & --- \\
Peak $\xir$ at 1 kHz            & 62.5             & 62.2 \\
Ceiling $\xir(0)$               & $62.4\!\pm\!1.2$ & $62.7\!\pm\!0.9$ \\
$\xi_\infty$                    & $21.7\!\pm\!0.9$ & $15.8\!\pm\!0.5$ \\
CMW at 1 kHz                    & 5.2              & 2.4 \\
$\CMW_0$ (fit)                  & 5.38             & 3.6$^{\ddagger}$ \\
$\fc$                           & 179 kHz          & 1.4 kHz \\
$\tau$                          & 0.90 $\mu$s      & 114 $\mu$s \\
Loss-peak corner                & 182 kHz          & --- \\
$\beta$ ($\alpha$)              & 1.00 ($-2.7^\circ$) & --- \\
$\Ecross$ (MV/cm)               & 0.221            & 0.081 \\
$R_s$ ($\Omega$)                & 267              & --- \\
$f_{RC}$                        & 1.81 MHz         & $\approx$220 kHz \\
$c_{IL}$ ($\mu$F/cm$^2$)        & 2.95             & 2.34 \\
$R_{IL}$ ($\Omega\cdot$cm$^2$)  & 0.31             & 48.8 \\
$\varepsilon_0/c_{IL}$ (nm)     & 0.301            & 0.379 \\
\midrule
\multicolumn{3}{@{}l}{\textit{Stack A, 300 K / 10 K}}\\
$2P_r$ ($\mu$C/cm$^2$)          & \multicolumn{2}{c}{33.1 / 20.3} \\
$E_c$ (MV/cm)                   & \multicolumn{2}{c}{0.93 / 1.17} \\
Imprint (V)                     & \multicolumn{2}{c}{$+0.017$ / $+0.009$} \\
CMW at 1 kHz                    & \multicolumn{2}{c}{5.2 / 2.2} \\
$\fc$                           & \multicolumn{2}{c}{179 kHz / $\lesssim$1 kHz} \\
\bottomrule
\end{tabular}
\\[2pt]
\raggedright\scriptsize $^{\dagger}$To be confirmed; used for field-unit
conversion. $^{\ddagger}$Extrapolated.
\end{table}

%=======================================================================
%  ACKNOWLEDGMENT
%=======================================================================
\section*{Acknowledgment}
The authors thank TSRI and the NTHU College of Semiconductor Research for
fabrication and measurement support.

%=======================================================================
%  REFERENCES
%=======================================================================
\begin{thebibliography}{99}
\itemsep 0pt

\bibitem{boscke2011}
T.~S. B\"oscke, J.~M\"uller, D.~Br\"auhaus, U.~Schr\"oder, and U.~B\"ottger,
``Ferroelectricity in hafnium oxide thin films,''
\emph{Appl. Phys. Lett.}, vol.~99, no.~10, 102903, 2011.

\bibitem{muller2012}
J.~M\"uller \emph{et al.},
``Ferroelectricity in simple binary ZrO$_2$ and HfO$_2$,''
\emph{Nano Lett.}, vol.~12, no.~8, pp.~4318--4323, 2012.

\bibitem{cho2024beol}
S.-H. Cho \emph{et al.},
``Exploring BEOL-compatible ferroelectricity in ultra-thin hafnium--zirconium
oxide,'' \emph{IEEE J. Electron Devices Soc.}, vol.~12, 2024.

\bibitem{schroeder2022}
U.~Schroeder, M.~H. Park, T.~Mikolajick, and C.~S. Hwang,
``The fundamentals and applications of ferroelectric HfO$_2$,''
\emph{Nat. Rev. Mater.}, vol.~7, pp.~653--669, 2022.

\bibitem{paul2025review}
A.~Paul, G.~Kumar, A.~Das, G.~Larrieu, and S.~De,
``Hafnium oxide-based ferroelectric field effect transistors: from materials
and reliability to applications in storage-class memory and in-memory
computing,'' \emph{J. Appl. Phys.}, vol.~138, no.~1, 010701, 2025.

\bibitem{sunil2024cam}
A.~Sunil \emph{et al.},
``Ferroelectric field effect transistors--based content-addressable
storage-class memory,'' \emph{Adv. Intell. Syst.}, 2024.

\bibitem{das2026cryo}
A.~Das \emph{et al.},
``Sub-2 nm equivalent-oxide-thickness ferroelectric transistors for cryogenic
memory and computing,'' \emph{ACS Nano}, 2026,
doi: 10.1021/acsnano.5c16255.

\bibitem{sk2025stdp}
S.~M. Sk \emph{et al.},
``Spike-timing dependent learning dynamics in silicon-doped hafnium-oxide
based ferroelectric field effect transistors,''
\emph{IEEE J. Electron Devices Soc.}, vol.~13, 2025.

\bibitem{mukherjee2023edl}
S.~Mukherjee \emph{et al.},
``Capacitive memory window with non-destructive read in ferroelectric
capacitors,'' \emph{IEEE Electron Device Lett.}, vol.~44, no.~7,
pp.~1092--1095, 2023.

\bibitem{mukherjee2023iedm}
S.~Mukherjee \emph{et al.},
in \emph{IEDM Tech. Dig.}, 2023.

\bibitem{mukherjee2024iedm}
S.~Mukherjee \emph{et al.},
in \emph{IEDM Tech. Dig.}, 2024.

\bibitem{lee2025aem}
M.~Lee \emph{et al.},
``Experimental demonstration of capacitive in-memory computing using
ferroelectric HZO memcapacitors,''
\emph{Adv. Electron. Mater.}, vol.~11, e00588, 2025.

\bibitem{mukherjee2025edl}
S.~Mukherjee \emph{et al.},
\emph{IEEE Electron Device Lett.}, vol.~46, no.~4, p.~576, 2025.

\bibitem{pesic2016}
M.~Pe\v{s}i\'c \emph{et al.},
``Physical mechanisms behind the field-cycling behavior of HfO$_2$-based
ferroelectric capacitors,'' \emph{Adv. Funct. Mater.}, vol.~26, no.~25,
pp.~4601--4612, 2016.

\bibitem{schenk2014}
T.~Schenk \emph{et al.},
``Electric field cycling behavior of ferroelectric hafnium oxide,''
\emph{ACS Appl. Mater. Interfaces}, vol.~6, no.~22, pp.~19744--19751, 2014.

\bibitem{grimley2016}
E.~D. Grimley \emph{et al.},
``Structural changes underlying field-cycling phenomena in ferroelectric
HfO$_2$ thin films,'' \emph{Adv. Electron. Mater.}, vol.~2, 1600173, 2016.

\bibitem{kim2023vacancy}
S.~Kim \emph{et al.},
``Role of oxygen vacancies in ferroelectric or resistive switching hafnium
oxide,'' \emph{Nano Convergence}, vol.~10, 55, 2023,
doi: 10.1186/s40580-023-00403-4.

\bibitem{segantini2023strain}
G.~Segantini \emph{et al.},
``Interplay between strain and defects at the interfaces of ultra-thin
\hzo{}-based ferroelectric capacitors,''
\emph{Adv. Electron. Mater.}, vol.~9, 2300171, 2023.

\bibitem{koroleva2022}
A.~Koroleva \emph{et al.},
``Retention improvement of HZO-based ferroelectric capacitors with TiO$_2$
insets,'' \emph{ACS Omega}, vol.~7, no.~51, pp.~47084--47093, 2022.

\bibitem{alcala2023}
R.~Alcala \emph{et al.},
``The electrode--ferroelectric interface as the primary constraint on
endurance and retention in HZO-based ferroelectric capacitors,''
\emph{Adv. Funct. Mater.}, vol.~33, no.~43, 2303261, 2023,
doi: 10.1002/adfm.202303261.

\bibitem{schenk2014apr}
T.~Schenk \emph{et al.},
``About the deformation of ferroelectric hystereses,''
\emph{Appl. Phys. Rev.}, vol.~1, 041103, 2014.

\bibitem{li2024imprint}
H.~Li \emph{et al.},
``Imprint in HfO$_2$ ferroelectric capacitors,''
in \emph{IEDM Tech. Dig.}, 2024.

\bibitem{cho2024wakeup}
S.-H. Cho \emph{et al.},
``Unraveling the wake-up mechanism in ultrathin ferroelectric \hzo{},''
\emph{IEEE Trans. Electron Devices}, vol.~71, 2024.

\bibitem{de2024tradeoff}
S.~De \emph{et al.},
``Trade-off between thermal budget and thickness scaling: a bottleneck on the
quest for BEOL-compatible ultra-thin ferroelectric films,''
in \emph{Proc. IEEE VLSI-TSA}, 2024.

\bibitem{iung2025}
J.~Iung \emph{et al.},
``Oxygen vacancy distribution and phase composition in scaled \hzo{}-based
ferroelectric capacitors,'' \emph{Appl. Phys. Lett.}, 2025.

\bibitem{kumar2025review}
G.~Kumar \emph{et al.},
``Insights into oxygen vacancy effects on ferroelectric behavior of hafnium
oxide: a review,'' \emph{Phys. Status Solidi A}, 2025.

\bibitem{rana2025cp}
S.~M. Rana \emph{et al.},
``Trapping dynamics and endurance in HfO$_2$-FeFETs: an insight from charge
pumping,'' \emph{IEEE Electron Device Lett.}, vol.~46, 2025.

\bibitem{raffel2025}
C.~Raffel \emph{et al.},
``Defect dynamics in silicon-doped HfO$_2$-based front-end-of-line FeFETs,''
\emph{IEEE Trans. Electron Devices}, 2025.

\bibitem{senapati2025}
A.~Senapati \emph{et al.},
``Long-term reliability of naturally aged hafnium oxide ferroelectric
transistors for energy-efficient embedded memory,''
\emph{Cell Rep. Phys. Sci.}, 2025.

\bibitem{das2026paradox}
A.~Das \emph{et al.},
``Endurance paradox in hafnium-oxide-based silicon-channel ferroelectric
transistors,'' \emph{ACS Appl. Mater. Interfaces}, 2026.

\bibitem{das2025sisc}
A.~Das \emph{et al.},
``Fundamental insights into trap accumulation and endurance degradation in
HfO$_2$ FeFETs,'' in \emph{Proc. SISC}, 2025.

\bibitem{paul2026wox}
A.~Paul \emph{et al.},
``WO$_x$--HZO--PEALD-TiN interface-engineered ferroelectric capacitors with
$>2\times10^{10}$ endurance,'' in \emph{Proc. IEEE VLSI-TSA}, 2026.

\bibitem{tagantsev2002}
A.~K. Tagantsev, I.~Stolichnov, N.~Setter, J.~S. Cross, and M.~Tsukada,
``Non-Kolmogorov-Avrami switching kinetics in ferroelectric thin films,''
\emph{Phys. Rev. B}, vol.~66, 214109, 2002.

\bibitem{segantini2022ti}
G.~Segantini \emph{et al.},
``Ferroelectricity improvement in ultra-thin \hzo{} capacitors by the
insertion of a Ti interfacial layer,''
\emph{Phys. Status Solidi RRL}, vol.~16, 2100583, 2022.

\bibitem{kim2023impedance}
T.~Kim, J.~A. del Alamo, and D.~A. Antoniadis,
``AC impedance characteristics of ferroelectric \hzo{}: from 1 kHz to
10 GHz,'' in \emph{IEDM Tech. Dig.}, 2023.

\bibitem{jonscher1977}
A.~K. Jonscher,
``The `universal' dielectric response,''
\emph{Nature}, vol.~267, pp.~673--679, 1977.

\bibitem{jonscher1983}
A.~K. Jonscher,
\emph{Dielectric Relaxation in Solids}.
London, U.K.: Chelsea Dielectrics Press, 1983.

\bibitem{nicollian1982}
E.~H. Nicollian and J.~R. Brews,
\emph{MOS Physics and Technology}.
New York, NY, USA: Wiley, 1982.

\bibitem{schroder2006}
D.~K. Schroder,
\emph{Semiconductor Material and Device Characterization}, 3rd~ed.
Hoboken, NJ, USA: Wiley, 2006.

\bibitem{kremer2003}
F.~Kremer and A.~Sch\"onhals, Eds.,
\emph{Broadband Dielectric Spectroscopy}.
Berlin, Germany: Springer, 2003.

\bibitem{colecole1941}
K.~S. Cole and R.~H. Cole,
``Dispersion and absorption in dielectrics. I. Alternating current
characteristics,'' \emph{J. Chem. Phys.}, vol.~9, no.~4, pp.~341--351, 1941.

\bibitem{li2021traps}
J.~Li, M.~Si, Y.~Qu, X.~Lyu, and P.~D. Ye,
``Quantitative characterization of ferroelectric/dielectric interface traps by
pulse measurements,'' \emph{IEEE Trans. Electron Devices}, vol.~68, no.~3,
pp.~1214--1220, 2021.

\bibitem{luo2020}
Y.-C. Luo, J.~Hur, P.~Wang, A.~I. Khan, and S.~Yu,
``Non-volatile, small-signal capacitance in ferroelectric capacitors,''
\emph{Appl. Phys. Lett.}, vol.~117, 073501, 2020.

\bibitem{fehlings2025}
L.~Fehlings \emph{et al.},
``Reliability of capacitive read in arrays of ferroelectric capacitors,''
in \emph{Proc. IEEE ISCAS}, 2025.

\end{thebibliography}

\end{document}
